\textbf{Pour cette première partie, aucune justification n'est demandée et une seule réponse est possible par question. Pour chaque question, reportez son numéro sur votre copie et indiquez votre réponse.}

Une réponse fausse ou l'absence de réponse n'enlève aucun point.

\subsubsection*{Question 1}
	La forme développée de l'expression \mbox{$(3x-2)^2$} est :
	\begin{enumerate}[label=\textbf{\alph*)}]
	\begin{multicols}{2}
		\item \mbox{$9x^2-4$}
		\item \mbox{$3x^2-12x+4$}
	\end{multicols}
	\begin{multicols}{2}
		\item \mbox{$9x^2-12x+4$}
		\item \mbox{$6x-4$}
	\end{multicols}
	\end{enumerate}

\subsubsection*{Question 2}
	\begin{minipage}[t]{0.5\linewidth}
	Dans le repère (O,I,J) ci-contre, la droite \mbox{$(\Delta)$} a pour équation :
	\begin{enumerate}[label=\textbf{\alph*)}]
	\begin{multicols}{2}
		\item \mbox{$y=2x+2$}
		\item \mbox{$y=-2x+2$}
	\end{multicols}
	\begin{multicols}{2}
		\item \mbox{$y=-x+2$}
		\item \mbox{$y=x+2$}
	\end{multicols}
	\end{enumerate}
	\end{minipage}\hfill
	\begin{minipage}[t]{0.45\linewidth}
	\begin{tikzpicture}[xscale=1,yscale=1,rotate=0,every node/.style={scale=1,rotate=0},baseline={(0,2.8)}]
		\clip(-1.6,-1.4)rectangle(3.6,3.2);
		\draw[blue!20, very thin](-1.6,-1.4)grid[xstep=0.2,ystep=0.2](3.6,3.2);
		\draw[blue!30,thick](-1.6,-1.4)grid(3.6,3.2);
		\draw[line width=0.3mm,->, >=triangle 45](-1.6,0)--(3.6,0);
		\draw[line width=0.3mm,->, >=triangle 45](0,-1.4)--(0,3.2);
		\draw[thick,shift={(0,0)}](2pt,0)--(-2pt,0) (0,2pt)--(0,-2pt) node[below left]{O};
		\draw[thick,shift={(1,0)}](2pt,0)--(-2pt,0) (0,2pt)--(0,-2pt) node[below]{I};
		\draw[thick,shift={(0,1)}](2pt,0)--(-2pt,0) (0,2pt)--(0,-2pt) node[left]{J};
		\draw[red, line width=0.3mm](-1.6,3.6)--(3.6,-1.6)node[pos=0.2,below left]{$(\Delta)$};
	\end{tikzpicture}
	\end{minipage}

\subsubsection*{Question 3}
	Dans une classe de première, 75\,\% des élèves étudient le grec.

	Les autres élèves étudient le latin : ils sont 9.

	Le nombre d'élèves de cette classe de première est égal à :
	\begin{enumerate}[label=\textbf{\alph*)}]
	\begin{multicols}{4}
		\item \mbox{$24$}
		\item \mbox{$30$}
		\item \mbox{$34$}
		\item \mbox{$36$}
	\end{multicols}
	\end{enumerate}

\subsubsection*{Question 4}
	Le prix d'un article augmente de 15\,\%.

	Cela signifie que le prix de l'article a été multiplié par :
	\begin{enumerate}[label=\textbf{\alph*)}]
	\begin{multicols}{4}
		\item \mbox{$\dfrac{15}{100}$}
		\item \mbox{$\np{1.15}$}
		\item \mbox{$\np{0.85}$}
		\item \mbox{$\np{1.115}$}
	\end{multicols}
	\end{enumerate}

\subsubsection*{Question 5}
	Parmi les réponses proposées, la valeur la plus proche de \mbox{$\dfrac{\np{150000}}{\np{3200}}$} est :
	\begin{enumerate}[label=\textbf{\alph*)}]
	\begin{multicols}{4}
		\item \mbox{$5$}
		\item \mbox{$50$}
		\item \mbox{$500$}
		\item \mbox{$\np{5000}$}
	\end{multicols}
	\end{enumerate}

\subsubsection*{Question 6}
	Une vidéo, d'une durée de 1 minute et 40 secondes, contient \np{2400}~images.

	Le nombre d'images par seconde est égal à :
	\begin{enumerate}[label=\textbf{\alph*)}]
	\begin{multicols}{2}
		\item\np{60} images/seconde
		\item\np{24} images/seconde
	\end{multicols}
	\begin{multicols}{2}
		\item\np{120} images/seconde
		\item\np{15} images/seconde
	\end{multicols}
	\end{enumerate}

\subsubsection*{Question 7}
	On considère une fonction $f$ définie sur $\mathbb{R}$ par :
	\[f(x)=\np{0.5}(x-3)^2+10\]

	On note \mbox{$\mathcal{C}$} sa courbe représentative dans un repère.

	Un seul des quatre points ci-dessous appartient à la courbe \mbox{$\mathcal{C}$}. Lequel ?
	\begin{enumerate}[label=\textbf{\alph*)}]
	\begin{multicols}{4}
		\item A~\mbox{$(-3;10)$}
		\item B~\mbox{$(3;\np{10.5})$}
		\item C~\mbox{$(3;10)$}
		\item D~\mbox{$(0;\np{19.5})$}
	\end{multicols}
	\end{enumerate}

\subsubsection*{Question 8}
	On considère le nombre \mbox{$A=\dfrac{10^{201}\times10^{-4}}{(10^2)^{100}}$}

	On peut affirmer que :
	\begin{enumerate}[label=\textbf{\alph*)}]
	\begin{multicols}{4}
		\item \mbox{$A=-\np{0.001}$}
		\item \mbox{$A=\np{0.0001}$}
		\item \mbox{$A=\np{0.001}$}
		\item \mbox{$A=\np{1000}$}
	\end{multicols}
	\end{enumerate}
